\documentclass[final-report]{../project-plan/report-template}

\usepackage{amsmath}
\usepackage{booktabs}
\usepackage{graphicx}
\hypersetup{hidelinks}

\university{Imperial College London}
\department{Department of Earth Science and Engineering}
\course{MSc in Applied Computational Science and Engineering}
\title{From Feedback to Action: Adaptive Support for Beginner Python Learning}
\author{Xiyi Xiong}
\email{}
\githubusername{}
\supervisors{Dr Marijan Beg}
\repository{https://github.com/ese-ada-lovelace-2025/irp-xx725}

\begin{document}

\maketitlepage

\section*{Abstract}

This report describes the design and evaluation of a notebook-native adaptive
support layer for beginner Python learning. On a frozen pack of 60 constructed
learner states, the constrained LLM selector obtained a mean blinded-judge score
of 3.741, compared with 3.729 for a fixed baseline and 3.883 for a rule-adaptive
baseline. All three paired 95\% bootstrap intervals included zero, so this
evaluation does not establish that the LLM selector outperformed either primary
baseline. The results concern action choice on the constructed benchmark and do
not measure real-student learning gains.

\section{Introduction}

% State the problem, motivation, research questions, and contributions. Update
% the project-plan framing to match the implemented system rather than intent.
Explain the gap between receiving feedback and choosing a useful next learning
action, then state the final research questions and contributions.

\section{Background and Related Work}

% Cover automated programming feedback, next-step hints, adaptive learning, and
% notebook-native educational tools. Add citations only from checked sources.
Position the implemented contribution against automated feedback and adaptive
support for introductory programming.

\section{System Design and Implementation}

% Describe evidence collection, deterministic gating, decision policies,
% transport, fallbacks, notebook presentation, and reproducibility boundaries.
Document the final architecture and the rationale for separating feedback,
support-format selection, and next-action selection.

\section{Evaluation Methodology}

The formal action-quality evaluation followed
\texttt{action-quality-protocol-v2} and
\texttt{evaluation-policy-suite-v3}. Protocol v2 superseded the unexecuted v1
suite after one v1 state was exposed during connectivity diagnosis. The v2 pack
assigned new identities to all states and replaced the complete affected
four-state cluster before any v2 model request. The resulting content-addressed
pack contained 60 blinded constructed states: 40 course-verified states, 10
generated-attempt states, and 10 generic-LLM states. It comprised 30
counterfactual pairs and 22 meaning-preserving invariance groups, while covering
the ten frozen course exercises.

Every state was evaluated under the three primary conditions
\texttt{fixed-v2}, \texttt{full-adaptive-v1}, and
\texttt{llm-next-step-v6}, plus the secondary \texttt{no-history-v1}
ablation. Each condition received the same converted decision input, and the
deterministic evidence gate was applied before policy or model selection. The
selector was \texttt{gpt-4o-mini-2024-07-18}; the judge was
\texttt{gpt-4.1-2025-04-14}. Both were accessed through OpenAI, but through
separate configuration namespaces.

The independent judge received a blinded state and one candidate action without
policy identity, selector model, fallback state, latency, reason codes, or other
candidates' decisions. Candidate starts were separated by at least 3000~ms.
Invalid output was repaired once; a second invalid output remained missing.
The analysis used seed \texttt{20260816}, Wilson intervals for rates, normal
intervals for mean scores, and 10,000 seeded paired bootstrap resamples for
primary-condition score differences. Executable hard constraints were analysed
separately from judge scores.

\section{Formal Results}

The selector run completed all 60 states in all four conditions. The judge
returned valid scores for 235 of 240 candidates. Table~\ref{tab:formal-scores}
shows the locked score summaries. The LLM selector's mean score was close to the
fixed baseline, while the rule-adaptive baseline had the highest descriptive
mean. The no-history ablation had the lowest descriptive mean.

\begin{table}[ht]
  \centering
  \caption{Blinded-judge action-quality scores on the constructed v2 state pack.}
  \label{tab:formal-scores}
  \begin{tabular}{lrrl}
    \toprule
    Condition & Judge $n$ & Mean & 95\% CI \\
    \midrule
    \texttt{fixed-v2} & 59 & 3.729 & [3.405, 4.053] \\
    \texttt{full-adaptive-v1} & 60 & 3.883 & [3.560, 4.206] \\
    \texttt{llm-next-step-v6} & 58 & 3.741 & [3.376, 4.107] \\
    \texttt{no-history-v1} & 58 & 3.207 & [2.804, 3.610] \\
    \bottomrule
  \end{tabular}
\end{table}

None of the three primary paired comparisons excluded zero. Full-adaptive minus
fixed was $+0.153$ ($n=59$, 95\% bootstrap CI $[-0.288, 0.576]$); LLM minus
fixed was $-0.035$ ($n=57$, CI $[-0.263, 0.211]$); and LLM minus full-adaptive
was $-0.138$ ($n=58$, CI $[-0.569, 0.293]$).

Executable diagnostics exposed trade-offs that the mean scores alone conceal.
The fixed condition had 0/60 hard-constraint violations, full-adaptive had 6/60
(10.0\%), and the LLM condition had 2/60 (3.3\%). The binary evidence-gate
classification field was correct for 60/60 records in every condition. Judge
coverage was 59/60 for fixed, 60/60 for full-adaptive, and 58/60 for both the LLM
and no-history conditions. Critical-error rates among completed judgements were
1/59, 1/60, 4/58, and 7/58 respectively. Invariance stability was 22/22 for
fixed, full-adaptive, and no-history, but 20/22 for the LLM condition. The LLM
selector used a rule fallback for 1/52 model-eligible states.

\section{Discussion}

The formal result does not support the hypothesis that the constrained LLM
selector produces higher-quality actions than either primary baseline on this
state pack. Its paired difference from fixed was near zero, and uncertainty in
all primary comparisons was wide enough to include advantages in either
direction. The rule-adaptive condition had the highest mean, but its six
executable progression violations show why a single subjective score should not
be treated as sufficient evidence of policy safety.

The secondary no-history condition's lower mean and higher critical-error rate
are consistent with attempt history being useful for action selection. This is
a mechanism-level observation rather than a causal learner-outcome claim: the
ablation was secondary, and the frozen primary paired analysis did not estimate
a no-history contrast. Similarly, the LLM condition's two hard violations and
lower invariance stability identify specific reliability targets, but v2 was
not tuned after observing them.

The fixed baseline combined a similar mean score to the LLM condition with no
hard violations and complete invariance stability. For this benchmark, the
extra selector complexity therefore did not yield a demonstrated action-quality
advantage. A future system version should address course-target validation and
then be evaluated on a new, uncontaminated state split rather than reusing v2.

\section{Limitations and Ethical Boundary}

No participants were recruited for the formal action-quality evaluation. The
formal runs measure configured systems on a constructed, blinded state pack;
they do not measure human learning, attainment, usability, confidence, or
behaviour. Consequently, this report must not claim real-student learning gains
or population-level learner benefit. Pilot ratings, development states, demo
traces, simulated outputs, and results observed during debugging are excluded
from formal results.

The state pack is constructed and limited to beginner-Python situations, so its
coverage and difficulty distribution need not represent authentic classroom
work. The selector and judge used different fixed model series but the same
provider, leaving provider-correlated judgement as a limitation. The judge
failed to return a valid grounded assessment for 5/240 candidates after the one
allowed repair, producing unequal condition coverage. Judge scores remain model
judgements rather than human educational judgements, and the reported intervals
describe this benchmark rather than a sampled learner population.

The frozen summary labels a 60-state binary classification measure as
``Needs-Evidence Accuracy''. Its denominator includes both evidence-insufficient
states and states correctly allowed to proceed; it must therefore be interpreted
as overall evidence-gate classification accuracy, not as a separately estimated
positive-class sensitivity. The LLM also fell back to a deterministic rule once
among 52 eligible states. Finally, the v1 exposure and v2 rebuild were recorded
before execution, but the transformation-based construction means that v1 and
v2 are not independent samples and must not be pooled.

\section{Reproducibility}

The formal run identifiers are \texttt{aq-v2-primary-001},
\texttt{aq-v2-judge-001}, and \texttt{aq-v2-summary-001}. The state-pack hash is
\texttt{A00F5F005903C18861A37FD92865B1DF589D8D1699232E6748A1370F8D230F96}.
The selector, judge, and statistics stages respectively recorded clean source
commits \texttt{d00dd256a096100297ffbddbdbbb5d9923009deb},
\texttt{c5c6e4d49e195677e6775e36ded2c0d3074f08a5}, and
\texttt{92f305d55187e781d09259eb2eabfed9b16b69b6}.

The locked statistics manifest links the source records, source manifest, judge
records, and judge manifest with SHA-256 values
\texttt{688A4B07E0777E16F3994914A2F817033EBD78F280761FF5416166C0D1B94809},
\texttt{FB81C4FC1B09DC9229AC7CBE53E422B4C10C667B9307EBDDFFE84FC0BF20FB66},
\texttt{AD439479E30AB395AECAB5E284F81E811ECCFA9B432863BDC2CD0BC3D00EACE6},
and \texttt{E0AF016C0DB8FB2FC2C50F0A9EFA7379AB117D00E331F801290CD8983CAB60EA}.
Reproduction consists of verifying these canonical hashes and regenerating the
JSON, CSV, and Markdown summary with seed \texttt{20260816} and 10,000
resamples. API keys, raw provider responses, and process-only proxy settings are
excluded from all manifests.

\section{Conclusion}

The project implemented a notebook-native, evidence-gated next-action system and
an auditable evaluation pipeline with frozen inputs, blinded judging,
write-once artefacts, and a reproducible hash chain. On the constructed v2
benchmark, the LLM selector did not demonstrate higher action quality than the
fixed or rule-adaptive baselines. The next empirical step is a newly frozen
evaluation of any reliability changes, followed---only with appropriate ethics
approval and consent---by research that measures outcomes with real learners.

% Add a checked bibliography when citations are introduced. The project-plan
% bibliography at ../project-plan/references.bib is a starting point, not an
% automatically authoritative final reference list.

\end{document}
